\documentclass[12pt,a4paper]{article}

% Packages
\usepackage[utf8]{inputenc}
\usepackage[T1]{fontenc}
\usepackage{geometry}
\usepackage{graphicx}
\usepackage{amsmath}
\usepackage{amssymb}
\usepackage{booktabs}
\usepackage{hyperref}
\usepackage{fancyhdr}
\usepackage{titlesec}
\usepackage{enumitem}
\usepackage{xcolor}
\usepackage{listings}
\usepackage{float}
\usepackage{caption}
\usepackage{subcaption}

\geometry{margin=2.5cm, top=3cm, bottom=2.5cm}

% Header with logo
\pagestyle{fancy}
\fancyhf{}
\fancyhead[L]{
\includegraphics[height=1.2cm]{uorlogo.jpg}}
\fancyhead[R]{\small Explainable AI --- MTech Data Science}
\fancyfoot[C]{\thepage}
\renewcommand{\headrulewidth}{0.4pt}
\renewcommand{\footrulewidth}{0.4pt}

% Code listing style
\lstset{
    basicstyle=\ttfamily\small,
    breaklines=true,
    frame=single,
    backgroundcolor=\color{gray!10},
    keywordstyle=\color{blue},
    commentstyle=\color{green!60!black},
    stringstyle=\color{red!70!black},
    language=Python
}

\begin{document}

% Title Page
\begin{titlepage}
    \centering
    
\includegraphics[width=4cm]{uorlogo.jpg}\\[1cm]
    {\Large \textbf{Nirma University}}\\[0.3cm]
    {\large Institute of Technology}\\[0.3cm]
    {\large Department of Computer Science \& Engineering}\\[1.5cm]
    {\LARGE \textbf{Explainable AI}}\\[0.5cm]
    {\Large \textbf{Lab 2: Explainability in Logistic Regression and SVM}}\\[2cm]
    \begin{tabular}{ll}
        \textbf{Student Name:} & Soham Dave \\[0.3cm]
        \textbf{Roll Number:} & 25MCD005 \\[0.3cm]
        \textbf{Program:} & MTech Data Science, Semester 2 \\[0.3cm]
        \textbf{Subject:} & Explainable AI \\[0.3cm]
        \textbf{Subject Guide:} & Dr. Nileshkumar Patel \\
    \end{tabular}
    \vfill
    {\large \today}
\end{titlepage}

\tableofcontents
\newpage

% -------------------------------------------------------
\section{Objective}
The objective of this lab is to demonstrate explainability in two widely-used classification models:
\begin{enumerate}
    \item \textbf{Logistic Regression} --- interpreting coefficients as log-odds contributions
    \item \textbf{Support Vector Machine (SVM)} --- visualizing the decision boundary, margin, and support vectors
\end{enumerate}

Both models are \textbf{intrinsically interpretable} when used with appropriate kernels and visualizations.

% -------------------------------------------------------
\section{Datasets}

\subsection{Synthetic Customer Churn Dataset}
\begin{itemize}
    \item \textbf{Task:} Binary classification (churn vs.\ no churn)
    \item \textbf{Samples:} 1,000 (generated using \texttt{make\_classification})
    \item \textbf{Features (5):} \texttt{feature\_1} through \texttt{feature\_5}, all informative with no redundant features
    \item \textbf{Target:} \texttt{churn} (0 = no churn, 1 = churn)
    \item \textbf{Generation Parameters:} \texttt{n\_informative=5}, \texttt{n\_redundant=0}, \texttt{random\_state=42}
\end{itemize}

\subsection{Iris Dataset (Binary Subset)}
\begin{itemize}
    \item \textbf{Task:} Binary classification (Setosa vs.\ Versicolor)
    \item \textbf{Source:} \texttt{sklearn.datasets.load\_iris}, filtered to classes 0 and 1
    \item \textbf{Samples:} 100 (50 per class)
    \item \textbf{Features Used (2):} Sepal length and petal length (selected for 2D visualization)
    \item \textbf{Target:} Species (0 = Setosa, 1 = Versicolor)
\end{itemize}

% -------------------------------------------------------
\section{Methodology}

\subsection{Part 1: Logistic Regression}

\subsubsection{Model Description}
Logistic Regression models the probability of the positive class using the logistic (sigmoid) function:
\begin{equation}
    P(y=1 \mid \mathbf{x}) = \sigma(\beta_0 + \beta_1 x_1 + \cdots + \beta_p x_p) = \frac{1}{1 + e^{-(\beta_0 + \boldsymbol{\beta}^T \mathbf{x})}}
\end{equation}

The coefficients $\beta_j$ represent the change in \textbf{log-odds} for a one-unit increase in feature $x_j$:
\begin{equation}
    \log\left(\frac{P(y=1)}{P(y=0)}\right) = \beta_0 + \sum_{j=1}^{p} \beta_j x_j
\end{equation}

\subsubsection{Implementation Steps}
\begin{enumerate}
    \item Generated synthetic binary classification data with 5 informative features
    \item Trained \texttt{LogisticRegression} using the \texttt{liblinear} solver
    \item Extracted and printed model coefficients and intercept
    \item Plotted absolute coefficient magnitudes as a bar chart to visualize feature importance
\end{enumerate}

\subsubsection{Explainability Analysis}

\begin{figure}[H]
    \centering
    \includegraphics[width=0.85\textwidth]{images_2/output_1.png}
    \caption{Logistic Regression coefficient magnitudes for the customer churn prediction model. Feature 2 has the strongest influence on churn prediction.}
    \label{fig:lr_coefficients}
\end{figure}

The trained model yielded the following coefficients:
\begin{center}
\begin{tabular}{lcc}
    \toprule
    \textbf{Feature} & \textbf{Coefficient} & \textbf{Effect on Churn} \\
    \midrule
    feature\_2 & $+0.8422$ & Increases churn probability \\
    feature\_1 & $-0.7461$ & Decreases churn probability \\
    feature\_5 & $-0.7187$ & Decreases churn probability \\
    feature\_3 & $+0.6936$ & Increases churn probability \\
    feature\_4 & $-0.3051$ & Decreases churn probability \\
    \midrule
    Intercept & $+0.7327$ & Baseline log-odds \\
    \bottomrule
\end{tabular}
\end{center}

\textbf{Interpretation:} A positive coefficient means that increasing the feature value increases the log-odds of churning. For example, a one-unit increase in \texttt{feature\_2} increases the log-odds of churn by 0.8422, which translates to an odds ratio of $e^{0.8422} \approx 2.32$ --- meaning the odds of churn more than double.

The positive intercept (0.7327) indicates that even without considering any features, there is an inherent bias toward churn in this dataset.

% -------------------------------------------------------
\subsection{Part 2: Support Vector Machine}

\subsubsection{Model Description}
A linear SVM finds the hyperplane that maximizes the margin between two classes:
\begin{equation}
    f(\mathbf{x}) = \mathbf{w}^T \mathbf{x} + b
\end{equation}

where $\mathbf{w}$ is the weight vector (normal to the hyperplane) and $b$ is the bias term. The decision boundary is defined by $f(\mathbf{x}) = 0$, and the margin boundaries are $f(\mathbf{x}) = \pm 1$.

The optimization objective is:
\begin{equation}
    \min_{\mathbf{w}, b} \frac{1}{2}\|\mathbf{w}\|^2 \quad \text{subject to} \quad y_i(\mathbf{w}^T \mathbf{x}_i + b) \geq 1, \quad \forall i
\end{equation}

\textbf{Support vectors} are the data points that lie exactly on the margin boundaries ($y_i f(\mathbf{x}_i) = 1$).

\subsubsection{Implementation Steps}
\begin{enumerate}
    \item Loaded the Iris dataset and filtered to two classes (Setosa and Versicolor)
    \item Selected two features (sepal length, petal length) for 2D visualization
    \item Trained \texttt{SVC(kernel='linear')} from scikit-learn
    \item Extracted the weight vector $\mathbf{w}$ and intercept $b$
    \item Computed the decision boundary line: $x_2 = -\frac{w_1}{w_2} x_1 - \frac{b}{w_2}$
    \item Computed margin lines at $\pm \frac{1}{w_2}$ offset from the decision boundary
    \item Plotted data points (colored by class), support vectors (circled), decision boundary, and margins
\end{enumerate}

\subsubsection{Explainability Analysis}

\begin{figure}[H]
    \centering
    \includegraphics[width=0.9\textwidth]{images_2/output_2.png}
    \caption{SVM decision boundary visualization. Data points are colored by class (Setosa vs.\ Versicolor). Support vectors (circled) define the margin. The dashed lines indicate the margin boundaries.}
    \label{fig:svm_boundary}
\end{figure}

\textbf{Interpretation:}
\begin{itemize}
    \item The \textbf{decision boundary} (solid line) separates Setosa from Versicolor flowers. Points above the line are classified as one class, points below as the other.
    \item The \textbf{support vectors} (circled points) are the critical samples that define the boundary. These are the ``hardest to classify'' examples --- removing any support vector would shift the boundary.
    \item The \textbf{margin width} indicates model confidence. A wider margin means the two classes are well-separated. In this binary Iris problem, the classes are highly separable, resulting in a clear wide margin.
    \item The \textbf{weight vector} $\mathbf{w}$ reveals that petal length has a larger weight than sepal length, indicating it is the more discriminative feature. This aligns with botanical knowledge --- Setosa flowers have distinctly shorter petals.
\end{itemize}

% -------------------------------------------------------
\section{Conclusion}
This lab demonstrated that both Logistic Regression and linear SVMs are inherently interpretable:
\begin{itemize}
    \item \textbf{Logistic Regression} coefficients directly quantify each feature's contribution to the prediction in terms of log-odds. This makes it easy to identify the most influential factors and their direction of effect.
    \item \textbf{Linear SVM} provides geometric interpretability through its decision boundary, support vectors, and margin. The weight vector reveals feature importance, and the support vectors identify the most informative training examples.
\end{itemize}

Both methods are widely used in domains requiring transparent decision-making, such as healthcare and finance, where practitioners need to understand and justify model predictions.

\end{document}
